\documentclass[11pt,letterpaper]{article}
\usepackage[T1]{fontenc}
\usepackage[utf8]{inputenc}
\usepackage{amsmath,amsthm,mathtools}
\usepackage{newtxtext,newtxmath}
\usepackage[margin=0.92in,headheight=14pt]{geometry}
\usepackage{microtype}
\usepackage{booktabs,array}
\usepackage{xcolor}
\usepackage{enumitem}
\usepackage{listings}
\usepackage{fancyhdr}
\usepackage{hyperref}
\usepackage{bookmark}
\definecolor{linkblue}{RGB}{28,65,104}
\hypersetup{colorlinks=true,linkcolor=linkblue,citecolor=linkblue,urlcolor=linkblue,
  pdftitle={A proof of the formulas in OEIS A260306},
  pdfauthor={Research note prepared with ChatGPT},
  pdfsubject={Ramanujan theta function, series inversion, Stirling numbers}}
\pagestyle{fancy}
\fancyhf{}
\fancyhead[L]{\small Ramanujan's theta function}
\fancyhead[R]{\small A proof of OEIS A260306}
\fancyfoot[C]{\thepage}
\renewcommand{\headrulewidth}{0.3pt}
\setlength{\parskip}{4pt}
\setlength{\parindent}{1.15em}
\setlength{\emergencystretch}{2em}
\allowdisplaybreaks[2]
\numberwithin{equation}{section}
\newtheorem{theorem}{Theorem}[section]
\newtheorem{lemma}[theorem]{Lemma}
\newtheorem{proposition}[theorem]{Proposition}
\newtheorem{corollary}[theorem]{Corollary}
\theoremstyle{remark}
\newtheorem{remark}[theorem]{Remark}
\newcommand{\coeff}[2]{[#1^{#2}]}
\newcommand{\Res}{\operatorname{Res}}
\newcommand{\dd}{\,\mathrm{d}}
\newcommand{\A}{\mathcal{A}}
\newcommand{\Sthree}[2]{\left\{\!\begin{matrix}#1\\#2\end{matrix}\!\right\}_{\geq 3}}
\lstset{language=Python,basicstyle=\ttfamily\small,columns=fullflexible,
  keepspaces=true,showstringspaces=false,breaklines=true,
  frame=single,framerule=0.3pt,rulecolor=\color{black!30},
  xleftmargin=0.3em,xrightmargin=0.3em,aboveskip=10pt,belowskip=8pt}
\title{An exact coefficient formula for\\[2pt]Ramanujan's theta function\\[6pt]
  \large A self-contained proof of both formulas in OEIS A260306}
\author{Research note}
\date{September 20, 2026}

\begin{document}
\maketitle
\thispagestyle{empty}
\begin{abstract}
For positive integers $N$, define $\theta(N)$ by
$\sum_{k=0}^{N-1}N^k/k!+\theta(N)N^N/N!=e^N/2$.
The numerators of the coefficients in its asymptotic expansion form OEIS
A260306. We prove both explicit formulas recorded in that entry: a double
sum involving Stirling numbers of the second kind and an equivalent triple
sum involving only binomial coefficients, factorials, and integer powers.
The proof first converts the defining expression into the difference of
two Laplace integrals. An analytic change of variable and a direct remainder
estimate establish the expansion to every order. Lagrange inversion then
reduces each coefficient to a single coefficient extraction. A finite
change-of-basis identity turns that extraction into the exact OEIS sums,
including all their signs and factorial factors. We also derive the
recurrence used in the entry, a one-sum formula using associated Stirling
numbers, and reproducible exact-arithmetic checks. The argument is a proof
for every nonnegative coefficient index, not an extrapolation from computed
initial values.
\end{abstract}

\section{The statement and its normalization}
\label{sec:statement}

For a positive integer $N$, let
\begin{equation}
 \sum_{k=0}^{N-1}\frac{N^k}{k!}
 +\theta(N)\frac{N^N}{N!}=\frac{e^N}{2}.
 \label{eq:def}
\end{equation}
We use $N$ for the large argument and $n$ for a coefficient index. Thus the
expansion under discussion is
\begin{equation}
 \theta(N)\sim\sum_{n\geq 0}\frac{c_n}{N^n}.
 \label{eq:asymp}
\end{equation}
In this article, \eqref{eq:asymp} means that, for each fixed integer $R\geq1$,
\begin{equation}
 \theta(N)=\sum_{n=0}^{R-1}\frac{c_n}{N^n}+O_R(N^{-R})
 \qquad(N\longrightarrow+\infty).
 \label{eq:poincare}
\end{equation}
It is not being used as an assertion that the infinite series converges.

Write $c_n=a(n)/b(n)$ in lowest terms, with $b(n)>0$. The OEIS entries
A260306 and A065973 specify $a(n)$ and $b(n)$, respectively
\cite{oeis260306,oeis065973}. Their initial coefficients are
\begin{equation}
 c_0=\frac13,\qquad c_1=\frac4{135},\qquad
 c_2=-\frac8{2835},\qquad c_3=-\frac{16}{8505},\qquad
 c_4=\frac{8992}{12629925}.
 \label{eq:initial}
\end{equation}
The formulas in the entry are formulas for the \emph{rational coefficients}
$c_n$, not for the integer numerators before reduction.

Let $S(L,j)$ denote a Stirling number of the second kind: the number of
partitions of an $L$-element set into $j$ nonempty blocks. In particular,
$S(0,0)=1$ and $S(L,0)=0$ for $L>0$.

\begin{theorem}[The two OEIS formulas]
\label{thm:main}
The expansion \eqref{eq:poincare} exists to every order. For every integer
$n\geq0$, its coefficient $c_n$ satisfies
\begin{equation}
\begin{split}
 c_n={}&2^n(3n+2)!
 \sum_{i=1}^{2n+1}\sum_{j=1}^{i}\\[-2pt]
 &\quad\times
 \frac{(-1)^{j+1}2^i S(2n+i+j+1,j)}
 {(2n+i+j+1)!\,(2n-i+1)!\,(i-j)!\,(n+i+1)}.
\end{split}
\label{eq:double}
\end{equation}
Equivalently,
\begin{equation}
\begin{split}
 c_n={}&\frac{2^n(3n+2)!}{(2n+1)!}
 \sum_{i=1}^{2n+1}\sum_{j=1}^{i}\sum_{k=1}^{j}\\[-2pt]
 &\quad\times
 \frac{(-1)^{k+1}2^i k^{2n+i+j+1}
       \binom{2n+1}{i}\binom{i}{j}\binom{j}{k}}
 {(2n+i+j+1)!\,(n+i+1)}.
\end{split}
\label{eq:triple}
\end{equation}
Consequently, the reduced numerators of either right-hand side are
A260306, and the reduced positive denominators are A065973.
\end{theorem}

The finite sums \eqref{eq:double}--\eqref{eq:triple} are the formulas displayed
in A260306 \cite{oeis260306}. O'Sullivan gives a broader study of Ramanujan's
approximation using saddle-point methods, including other closed forms
for its coefficients \cite{osullivan}. Our purpose here is to give a direct,
self-contained derivation of these \emph{particular} OEIS expressions.
No claim of novelty for Ramanujan's expansion or its general coefficient
theory is needed.

The decisive intermediate identity will be
\begin{equation}
 \boxed{\displaystyle
 c_n=-2^n n!\,[z^{2n+1}]
 \left(\frac{z^2}{2(e^z-1-z)}\right)^{n+1}.}
 \label{eq:key}
\end{equation}
The quotient in parentheses has a removable singularity at zero and
constant term $1$. Sections~\ref{sec:integral}--\ref{sec:lagrange} establish
\eqref{eq:key} analytically; Sections~\ref{sec:finite}--\ref{sec:triple}
convert it into the advertised sums by finite algebra.

\section{An exact integral representation}
\label{sec:integral}

For $x>0$ and $u>0$, define
\[
 \Gamma(x)=\int_0^\infty t^{x-1}e^{-t}\dd t,
 \qquad
 \Gamma(x,u)=\int_u^\infty t^{x-1}e^{-t}\dd t.
\]
Integration by parts, starting with $\Gamma(1,u)=e^{-u}$, gives
\[
 \Gamma(N,u)=(N-1)!e^{-u}\sum_{k=0}^{N-1}\frac{u^k}{k!}
 \qquad(N\in\mathbb Z_{\geq1}).
\]
Setting $u=N$ in this identity and rearranging \eqref{eq:def} yields
\begin{equation}
 \theta(N)=\frac{Ne^N}{N^N}
 \left(\frac12\Gamma(N)-\Gamma(N,N)\right).
 \label{eq:gamma}
\end{equation}
The right-hand side also defines a natural extension to every real $N>0$.
We use that extension in the analysis below; its restriction to positive
integers is exactly the original function.

In both gamma integrals substitute $t=Ne^x$. Since
\[
 t^{N-1}e^{-t}\dd t
 =N^Ne^{-N}\exp\{-N(e^x-1-x)\}\dd x,
\]
formula \eqref{eq:gamma} becomes
\begin{equation}
 \theta(N)=\frac N2\left(
   \int_{-\infty}^{0}e^{-NF(x)}\dd x
   -\int_0^\infty e^{-NF(x)}\dd x\right),
 \qquad F(x)=e^x-1-x.
 \label{eq:laplace}
\end{equation}
This equality is exact. The common exponential scale and all gamma factors
have disappeared. In particular, no separate use of Stirling's expansion
for $N!$ is required.

The function $F$ is positive away from zero, with
\[
 F(0)=F'(0)=0,\qquad F''(0)=1,\qquad F'(x)=e^x-1.
\]
Thus each integral is concentrated near the same nondegenerate minimum.
Their difference will cancel the terms that would otherwise produce
half-integer powers of $N^{-1}$.

\section{The inverse change of variable and the remainder}
\label{sec:analytic}

Define the entire function with a removable value at zero
\begin{equation}
 A(z)=\frac{2(e^z-1-z)}{z^2}
 =\sum_{r=0}^\infty\frac{2z^r}{(r+2)!}
 =1+\frac z3+\frac{z^2}{12}+\frac{z^3}{60}+\cdots.
 \label{eq:A}
\end{equation}
Near $z=0$, choose the analytic square root of $A(z)$ with value $1$.
The function $w=z\sqrt{A(z)}$ then has derivative $1$ at zero, so it has a
unique local analytic inverse
\begin{equation}
 z=\psi(w)=\sum_{m=1}^\infty b_mw^m,
 \qquad b_1=1,
 \qquad F(\psi(w))=\frac{w^2}{2}.
 \label{eq:inverse}
\end{equation}

On the real line the same change of variable can be written globally as
\[
 w=\operatorname{sgn}(x)\sqrt{2F(x)}.
\]
It is strictly increasing: away from zero its derivative is $F'(x)/w>0$,
and at zero its derivative is $1$. Its limits at the two ends of the real
line are $-\infty$ and $+\infty$, respectively. It therefore has a global
real inverse, agreeing with the analytic inverse $\psi$ near zero.
Differentiating \eqref{eq:inverse} gives
\begin{equation}
 \psi'(w)=\frac{w}{e^{\psi(w)}-1}\qquad(w\ne0).
 \label{eq:derivative}
\end{equation}

Substitution into \eqref{eq:laplace}, followed by reflection of the negative
half-axis, gives the exact identity
\begin{equation}
 \theta(N)=\frac N2\int_0^\infty e^{-Nw^2/2}
              \bigl(\psi'(-w)-\psi'(w)\bigr)\dd w.
 \label{eq:gaussian}
\end{equation}

\begin{proposition}[All-orders expansion]
\label{prop:asymp}
For every fixed integer $R\geq1$, \eqref{eq:poincare} holds with
\begin{equation}
 c_n=-2^{n+1}(n+1)!\,b_{2n+2},\qquad n\geq0.
 \label{eq:bcoeff}
\end{equation}
\end{proposition}

\begin{proof}
Choose $\delta>0$ small enough that the power series for $\psi$ is analytic
on a neighborhood of $[-\delta,\delta]$. Its derivative difference has the
odd Taylor expansion
\begin{equation}
 \psi'(-w)-\psi'(w)
 =-2\sum_{n=0}^{R-1}(2n+2)b_{2n+2}w^{2n+1}
   +O_R(w^{2R+1})
 \quad(0\leq w\leq\delta).
 \label{eq:oddpart}
\end{equation}
The required Gaussian moments are elementary:
\begin{equation}
 \int_0^\infty e^{-Nw^2/2}w^{2n+1}\dd w
 =\frac{2^n n!}{N^{n+1}}.
 \label{eq:moments}
\end{equation}
Indeed, substitute $v=Nw^2/2$.

We now justify replacing the local expansion by these moments. For
$w\geq\delta$, the negative inverse branch satisfies
\[
 0<\psi'(-w)
 =\frac{w}{1-e^{\psi(-w)}}
 \leq\frac{w}{1-e^{\psi(-\delta)}}.
\]
For the positive branch, $e^{\psi(w)}-1=\psi(w)+w^2/2\geq w^2/2$, so
\[
 0<\psi'(w)\leq\frac2w\leq\frac{2w}{\delta^2}.
\]
Consequently,
\[
 \frac N2\int_\delta^\infty e^{-Nw^2/2}
       |\psi'(-w)-\psi'(w)|\dd w
 \leq C_\delta N\int_\delta^\infty we^{-Nw^2/2}\dd w
 =C_\delta e^{-N\delta^2/2}.
\]
Thus the nonlocal part is exponentially small.

On $[0,\delta]$, the remainder in \eqref{eq:oddpart} contributes at most
\[
 C_R N\int_0^\delta e^{-Nw^2/2}w^{2R+1}\dd w
 \leq C_R\,2^R R!\,N^{-R}.
\]
For each of the finitely many polynomial terms, extending its integral
from $\delta$ to infinity changes the answer by an exponentially small
quantity. One can see this directly by substituting $v=w^2/2$ and
integrating $e^{-Nv}v^n$ by parts a finite number of times.

Finally, inserting \eqref{eq:oddpart} and \eqref{eq:moments} into
\eqref{eq:gaussian} gives
\[
 \frac N2\cdot(-2)(2n+2)b_{2n+2}
        \frac{2^n n!}{N^{n+1}}
 =\frac{-2^{n+1}(n+1)!b_{2n+2}}{N^n}.
\]
The stated coefficients and remainder follow.
\end{proof}

The parity cancellation in \eqref{eq:oddpart} explains why only integer
powers $N^{-n}$ occur. The argument proves the expansion for real
$N\to+\infty$, hence in particular along the positive integers.
Its coefficients are unique: if two expansions first differ at index
$n$, subtracting them and multiplying by $N^n$ gives a nonzero constant
that would have to tend to zero.

\section{Lagrange inversion and the coefficient identity}
\label{sec:lagrange}

We include the particular form of Lagrange inversion that is needed, so
that no coefficient-extraction step is left implicit.

\begin{lemma}[Inverse-series coefficient formula]
\label{lem:lagrange}
Suppose $w=z\phi(z)$, where $\phi$ is analytic near zero and $\phi(0)=1$.
Let $z=g(w)$ be its local inverse. For every $m\geq1$,
\begin{equation}
 [w^m]g(w)=\frac1m[z^{m-1}]\phi(z)^{-m}.
 \label{eq:lagrange}
\end{equation}
\end{lemma}

\begin{proof}
Differentiation and a change of variable in the residue give
\begin{align*}
 m[w^m]g(w)
 &=\Res_{w=0}\frac{g'(w)}{w^m}\dd w\\
 &=\Res_{z=0}\frac{1}{(z\phi(z))^m}\dd z
 =[z^{m-1}]\phi(z)^{-m}.
\end{align*}
All residues can be taken on sufficiently small contours on which the
analytic change of variable is invertible.
\end{proof}

Apply the lemma with $\phi(z)=\sqrt{A(z)}$ and $m=2n+2$. Then
\[
 b_{2n+2}
 =\frac1{2n+2}[z^{2n+1}]A(z)^{-n-1}.
\]
Together with \eqref{eq:bcoeff}, this proves the key identity
\eqref{eq:key}. In particular, every $c_n$ is rational, since the
power series of $A$ has rational coefficients.

\begin{corollary}[Residue form]
\label{cor:residue}
For every $n\geq0$,
\begin{equation}
 c_n=-\frac{n!}{2}\Res_{z=0}
             \frac{\dd z}{(e^z-1-z)^{n+1}}.
 \label{eq:residue}
\end{equation}
\end{corollary}
\begin{proof}
In \eqref{eq:key}, replace $A(z)^{-n-1}$ by
$z^{2n+2}/\bigl(2^{n+1}F(z)^{n+1}\bigr)$. The coefficient of $z^{2n+1}$
is $2^{-n-1}$ times the coefficient of $z^{-1}$ in $F(z)^{-n-1}$.
\end{proof}

Formula \eqref{eq:key} is already a compact exact answer. The remaining
work is to show why its finite expansion has precisely the less obvious
shape recorded in OEIS.

\section{A finite change of basis for negative powers}
\label{sec:finite}

The factorial $(3n+2)!$ and the factor $n+i+1$ in the OEIS formulas come
from a general finite algebraic identity.

\begin{lemma}[Finite negative-power identity]
\label{lem:finite}
Let $p\geq1$ and $M\geq1$ be integers, and let $B(z)$ be a formal power
series over $\mathbb Q$ or $\mathbb C$ with $B(0)=1$. Then
\begin{equation}
 [z^M]B(z)^{-p}
 =\frac{(p+M)!}{(p-1)!}
   \sum_{i=1}^{M}
   \frac{(-1)^i[z^M]B(z)^i}{i!\,(M-i)!\,(p+i)}.
 \label{eq:finite}
\end{equation}
\end{lemma}

\begin{proof}
Since $B(z)-1$ has zero constant term, only powers at most $M$ can affect
a coefficient of degree $M$. Therefore
\begin{equation}
 B(z)^{-p}\equiv
 \sum_{r=0}^{M}(-1)^r\binom{p+r-1}{r}(B(z)-1)^r
 \pmod{z^{M+1}}.
 \label{eq:truncated}
\end{equation}
Expand $(B-1)^r$ in powers of $B$ and interchange the two finite sums.
The coefficient multiplying $B^i$ is
\begin{align*}
 \sum_{r=i}^{M}(-1)^r\binom{p+r-1}{r}
                     \binom ri(-1)^{r-i}
 &=(-1)^i\sum_{r=i}^{M}\binom{p+r-1}{r}\binom ri\\
 &=(-1)^i\binom{p+i-1}{i}
                \sum_{r=i}^{M}\binom{p+r-1}{r-i}\\
 &=(-1)^i\binom{p+i-1}{i}\binom{p+M}{M-i}\\
 &=\frac{(-1)^i(p+M)!}
        {(p-1)!\,i!\,(M-i)!\,(p+i)}.
\end{align*}
The penultimate equality is the hockey-stick identity, obtained by repeated
application of Pascal's identity. The term $i=0$ multiplies $B^0=1$ and has
zero coefficient of $z^M$, because $M\geq1$. Taking $[z^M]$ proves the lemma.
\end{proof}

\begin{remark}[An integral interpretation of the denominator]
The polynomial in $X$ produced by the proof can also be written as
\begin{equation}
 T_{p,M}(X)=\frac{(p+M)!}{(p-1)!M!}
        \int_0^1 t^{p-1}(1-Xt)^M\dd t.
 \label{eq:beta}
\end{equation}
Expanding $(1-Xt)^M$ shows that the coefficient of $X^i$ includes
$\int_0^1t^{p+i-1}\dd t=1/(p+i)$. Thus the denominator in
\eqref{eq:finite} has a simple finite-polynomial origin.
\end{remark}

Now put
\begin{equation}
 M=2n+1,\qquad p=n+1,\qquad B(z)=A(z).
 \label{eq:specialization}
\end{equation}
Multiplying \eqref{eq:finite} by $-2^n n!$, as prescribed by
\eqref{eq:key}, gives
\begin{equation}
 c_n=2^n(3n+2)!
 \sum_{i=1}^{M}
 \frac{(-1)^{i+1}[z^M]A(z)^i}
      {i!\,(M-i)!\,(n+i+1)}.
 \label{eq:intermediate}
\end{equation}
Here $(p+M)!=(3n+2)!$ and $(p-1)!=n!$, which cancels the $n!$
from \eqref{eq:key}. These are exactly the factorials needed in the target.

\section{Extraction with Stirling numbers: the double sum}
\label{sec:double}

The exponential generating function for Stirling numbers is
\begin{equation}
 (e^z-1)^j=j!\sum_{L\geq0}S(L,j)\frac{z^L}{L!}.
 \label{eq:stirling-egf}
\end{equation}
For completeness, expanding $(e^z-1)^j$ gives, for $L,j\geq1$,
\begin{equation}
 j!S(L,j)=\sum_{k=0}^{j}(-1)^{j-k}\binom jk k^L.
 \label{eq:finite-difference}
\end{equation}
The right-hand side counts surjections from an $L$-element set to a
$j$-element set by inclusion--exclusion; dividing by $j!$ forgets the
labels of the nonempty fibers. This proves \eqref{eq:stirling-egf} and
\eqref{eq:finite-difference} directly.

From \eqref{eq:A}, for $i\geq1$,
\begin{align}
 [z^M]A(z)^i
 &=2^i[z^{M+2i}](e^z-1-z)^i \notag\\
 &=2^i\sum_{j=0}^{i}(-1)^{i-j}\binom ij
          [z^{M+i+j}](e^z-1)^j \notag\\
 &=2^i\sum_{j=1}^{i}(-1)^{i-j}\binom ij
          \frac{j!S(M+i+j,j)}{(M+i+j)!}.
 \label{eq:Apositive}
\end{align}
The second line follows by writing $e^z-1-z=(e^z-1)-z$.
In the third line, the $j=0$ term vanishes because it asks for a positive
coefficient of the constant series $1$.

Insert \eqref{eq:Apositive} into \eqref{eq:intermediate}. Two elementary
simplifications are responsible for the final form:
\[
 (-1)^{i+1}(-1)^{i-j}=(-1)^{j+1},
 \qquad
 \frac1{i!}\binom ij j!=\frac1{(i-j)!}.
\]
It follows that
\begin{equation}
 c_n=2^n(3n+2)!
 \sum_{i=1}^{M}\sum_{j=1}^{i}
 \frac{(-1)^{j+1}2^i S(M+i+j,j)}
      {(M+i+j)!\,(M-i)!\,(i-j)!\,(n+i+1)}.
 \label{eq:double-M}
\end{equation}
Substituting $M=2n+1$ is exactly \eqref{eq:double}. This proves the
Stirling-number formula for every $n\geq0$.

All bounds and factorials are valid at the endpoints: $M\geq1$,
$1\leq j\leq i\leq M$, and therefore $M-i$ and $i-j$ are nonnegative.
No zero factorial is omitted, and $n+i+1$ is strictly positive.

\section{From the double sum to the triple sum}
\label{sec:triple}

Let $L=M+i+j$. This exponent is positive, so the $k=0$ term in
\eqref{eq:finite-difference} is zero. Multiplying that identity by
$(-1)^{j+1}$ gives
\begin{equation}
 (-1)^{j+1}j!S(L,j)
 =\sum_{k=1}^{j}(-1)^{k+1}\binom jk k^L.
 \label{eq:signidentity}
\end{equation}
Indeed, $(-1)^{j+1}(-1)^{j-k}=(-1)^{k+1}$.

Also,
\begin{equation}
 \frac{S(L,j)}{(M-i)!\,(i-j)!}
 =\frac1{M!}\binom Mi\binom ij\,j!S(L,j).
 \label{eq:binomials}
\end{equation}
Substituting \eqref{eq:binomials} and \eqref{eq:signidentity} into
\eqref{eq:double-M} yields
\begin{equation}
\begin{split}
 c_n={}&\frac{2^n(3n+2)!}{M!}
 \sum_{i=1}^{M}\sum_{j=1}^{i}\sum_{k=1}^{j}\\[-2pt]
 &\qquad\times
 \frac{(-1)^{k+1}2^i k^{M+i+j}
       \binom Mi\binom ij\binom jk}
      {(M+i+j)!\,(n+i+1)}.
\end{split}
\label{eq:triple-M}
\end{equation}
Replacing $M$ by $2n+1$ proves \eqref{eq:triple} and completes the proof of
Theorem~\ref{thm:main}. Since every exponent $L$ here is positive, no
convention concerning $0^0$ is involved.

The roles of the three summation indices are now explicit. The index $i$
arises from replacing a negative power of $A$ by a finite combination of
positive powers. The index $j$ comes from expanding
$((e^z-1)-z)^i$. The index $k$ is the inclusion--exclusion index that
eliminates the Stirling number. In particular, equivalence of the double
and triple sums alone would not identify them with Ramanujan's
coefficients; the integral and inversion arguments supply that essential
identification.

\section{The recurrence in the OEIS entry}
\label{sec:recurrence}

A recurrence provides a compact way to compute a long initial segment.
It is also useful to explain the separate treatment of the constant term
in the OEIS program \cite{oeis260306}.

Define
\begin{equation}
 u(w)=e^{\psi(w)}-1.
 \label{eq:u}
\end{equation}
The inverse equation $e^{\psi(w)}-1-\psi(w)=w^2/2$ implies the especially
simple relation
\begin{equation}
 u(w)=\psi(w)+\frac{w^2}{2}.
 \label{eq:upsi}
\end{equation}
Equivalently, $u-\log(1+u)=w^2/2$ near zero. Let
\begin{equation}
 H(w)=u'(w)=\sum_{m\geq0}h_mw^m,\qquad h_0=1.
 \label{eq:H}
\end{equation}
Differentiating the latter inverse equation gives
\begin{equation}
 u(w)H(w)=w(1+u(w)).
 \label{eq:quadratic}
\end{equation}
Since $u(w)=\sum_{j\geq0}h_jw^{j+1}/(j+1)$, comparison of coefficients of
$w^{m+1}$ gives, for $m\geq1$,
\[
 \sum_{j=0}^{m}\frac{h_jh_{m-j}}{j+1}=\frac{h_{m-1}}m.
\]
Separating the terms $j=0$ and $j=m$ yields
\begin{equation}
 \boxed{\displaystyle
 h_m=\frac{\displaystyle\frac{h_{m-1}}m
          -\sum_{j=1}^{m-1}\frac{h_{m-j}h_j}{j+1}}
         {\displaystyle1+\frac1{m+1}},\qquad h_0=1.}
 \label{eq:hrec}
\end{equation}
Its beginning is
\[
 H(w)=1+\frac23w+\frac1{12}w^2-\frac2{135}w^3
          +\frac1{864}w^4+\frac1{2835}w^5+\cdots.
\]

By \eqref{eq:upsi}, $H(w)=\psi'(w)+w$, hence
\[
 h_{2n+1}=(2n+2)b_{2n+2}+\delta_{n,0}.
\]
Formula \eqref{eq:bcoeff} is therefore equivalent to
\begin{equation}
 \boxed{\displaystyle
 c_n=\delta_{n,0}-2^n n!\,h_{2n+1}.}
 \label{eq:ch}
\end{equation}
For $n=0$, this says $c_0=1-h_1=1/3$. For $n\geq1$, it says
$c_n=-(2n)!!\,h_{2n+1}$, since $(2n)!!=2^nn!$.
These are precisely the recurrence and the constant-term exception
underlying the Maple code in A260306. Forgetting the Kronecker delta in
\eqref{eq:ch} would give $-2/3$ instead of $1/3$ at index zero.

\subsection{A second coefficient algorithm}

There is a separate exact computation directly from \eqref{eq:key}.
For fixed $p\geq1$, write
\[
 A(z)^{-p}=\sum_{m\geq0}f_mz^m,\qquad f_0=1,
 \qquad a_k=[z^k]A(z)=\frac2{(k+2)!}.
\]
The differential identity $A(A^{-p})'=-pA'A^{-p}$ gives
\begin{equation}
 f_m=-\frac1m\sum_{k=1}^{m}
              \bigl(m+(p-1)k\bigr)a_kf_{m-k}
 \qquad(m\geq1).
 \label{eq:powerrec}
\end{equation}
Use $p=n+1$ and $m=2n+1$, then multiply by $-2^nn!$.
This algorithm is independently implemented in the verification code,
alongside \eqref{eq:hrec}, both OEIS sums, and the next corollary.

To obtain all coefficients through $c_K$, recurrence \eqref{eq:hrec}
uses $O(K^2)$ rational arithmetic operations. A single double-sum
coefficient has $O(n^2)$ summands after its Stirling table is available;
the literal triple sum has $O(n^3)$ summands. These are arithmetic-operation
counts, not bit-complexity bounds. Numerators and denominators grow, and
alternating sums can involve severe cancellation, so the exact algorithms
use integers and rational numbers rather than floating-point arithmetic.

\section{A one-sum corollary with associated Stirling numbers}
\label{sec:associated}

Let $\Sthree{L}{k}$ count partitions of an $L$-element set into $k$ blocks,
each of size at least three. The associated exponential generating
function is
\begin{equation}
 \frac{(e^z-1-z-z^2/2)^k}{k!}
 =\sum_{L\geq0}\Sthree{L}{k}\frac{z^L}{L!}.
 \label{eq:associated-egf}
\end{equation}
This follows by forming $k$ labeled blocks of sizes at least three and
then dividing by $k!$ to remove their order. O'Sullivan's treatment also
relates the asymptotic coefficients to associated Stirling numbers
\cite{osullivan}; here the needed identity follows immediately from
\eqref{eq:key}.

\begin{corollary}
\label{cor:associated}
For every $n\geq0$,
\begin{equation}
 c_n=\sum_{k=1}^{2n+1}
       \frac{(-1)^{k+1}2^{n+k}(n+k)!}
            {(2n+2k+1)!}
       \Sthree{2n+2k+1}{k}.
 \label{eq:associated}
\end{equation}
\end{corollary}

\begin{proof}
Put $M=2n+1$. Instead of changing basis as in Lemma~\ref{lem:finite}, expand
$A^{-n-1}$ directly in powers of $A-1$:
\[
 [z^M]A(z)^{-n-1}
 =\sum_{k=1}^{M}(-1)^k\binom{n+k}{k}[z^M](A(z)-1)^k.
\]
Since $A(z)-1=2(e^z-1-z-z^2/2)/z^2$, \eqref{eq:associated-egf} gives
\[
 [z^M](A(z)-1)^k
 =\frac{2^k k!}{(M+2k)!}\Sthree{M+2k}{k}.
\]
Multiplication by $-2^nn!$ proves the result. The upper bound $k\leq M$
also follows combinatorially: $M+2k$ elements cannot form $k$ blocks of
size at least three unless $M+2k\geq3k$.
\end{proof}

For exact computation, these associated numbers satisfy, for $L,k\geq1$,
\begin{equation}
 \Sthree{L}{k}
 =k\Sthree{L-1}{k}
   +\binom{L-1}{2}\Sthree{L-3}{k-1}.
 \label{eq:associated-rec}
\end{equation}
Use the usual zero conventions outside the allowed range and
$\Sthree{0}{0}=1$. To prove the recurrence, distinguish one element. Either
its block has at least four elements, in which case removing it leaves
one of $k$ eligible blocks, or its block has exactly three elements, in
which case choose its two companions and partition the remaining $L-3$
elements. This supplies another integer-based check of \eqref{eq:associated}.

\section{Initial coefficients and worked cases}
\label{sec:examples}

The first case is visible without any summation. From \eqref{eq:A},
\[
 [z]A(z)^{-1}=-\frac13,
 \qquad
 c_0=-[z]A(z)^{-1}=\frac13.
\]
The double sum also has only one term when $n=0$: $i=j=1$, giving
\[
 2!\,\frac{2S(3,1)}{3!\,0!\,0!\,2}=\frac13.
\]
Thus the displayed OEIS sum itself needs no exceptional rule at $n=0$.

For $n=1$, put $a_1=1/3$, $a_2=1/12$, and $a_3=1/60$. The binomial
expansion of $A^{-2}$ yields
\begin{align*}
 [z^3]A(z)^{-2}
 &=-2a_3+6a_1a_2-4a_1^3\\
 &=-\frac1{30}+\frac16-\frac4{27}
 =-\frac2{135}.
\end{align*}
Therefore $c_1=-2[z^3]A(z)^{-2}=4/135$, confirming both the factor and the
sign in the general formula.

\begin{table}[htbp]
\centering
\small
\begin{tabular}{r r r}
\toprule
$n$ & Reduced numerator $a(n)$ & Positive denominator $b(n)$\\
\midrule
0 & 1 & 3\\
1 & 4 & 135\\
2 & $-8$ & 2835\\
3 & $-16$ & 8505\\
4 & 8992 & 12629925\\
5 & 334144 & 492567075\\
6 & $-698752$ & 1477701225\\
7 & $-23349012224$ & 39565450299375\\
8 & 1357305243136 & 2255230667064375\\
9 & 6319924923392 & 6765692001193125\\
10 & $-8773495082018816$ & 7002491221234884375\\
\bottomrule
\end{tabular}
\caption{Coefficients generated by exact arithmetic from the proved
recurrence. The displayed values agree with the corresponding prefixes
of A260306 and A065973 \cite{oeis260306,oeis065973}.}
\label{tab:coefficients}
\end{table}

The table illustrates why it is important to reduce each rational
coefficient before extracting its numerator or denominator. Neither the
individual summands nor the unreduced factorial prefactor alone determines
the two OEIS sequences.

\section{Reproducible computation and a numerical illustration}
\label{sec:verification}

The accompanying Python code uses the standard-library
\texttt{fractions.Fraction} class for every exact coefficient computation.
The executed checks are recorded in \texttt{artifacts/checks.json}.
Specifically, the inverse-function recurrence generated $c_0$ through
$c_{117}$. For $0\leq n\leq20$, these coefficients were compared exactly
with the inverse-power algorithm, the double sum, and the associated
Stirling-number sum. For $0\leq n\leq10$, they were also compared with the
literal triple sum. Every comparison passed.

An external-data check compared the first 17 computed numerators and the
first 14 computed denominators with prefixes transcribed from the
consulted OEIS entries. The precise prefixes, source addresses, and access
date are in \texttt{artifacts/reference\_prefix.json}. All of those
comparisons passed as well. The 118 generated coefficients should not be
confused with 118 independently checked OEIS values: the external
comparison in this run covers only the stated prefixes. The optional
live b-file comparison is available in the script but was not completed
in the restricted execution environment.

To rerun the exact checks from the archive's root directory, use
\begin{lstlisting}[language=bash]
python3 code/verify.py
\end{lstlisting}
No third-party package and no network access is needed for that command.
The output includes \texttt{coefficients.csv}, containing separately the
index, reduced numerator, and reduced positive denominator. These checks
are finite consistency tests; the all-index assertion is established by
Theorem~\ref{thm:main}.

\subsection{Checking the asymptotic scale}

A separate script evaluates $\theta(N)$ in two ways: from the scaled
finite-sum definition and from the incomplete-gamma expression
\eqref{eq:gamma}. In the executed run, \texttt{mpmath} version 1.3.0 was used
at 100 decimal digits of working precision. The two evaluations differed
by less than $7.18\times10^{-99}$ for all tested
$N\in\{5,10,20,50,100,200,500,1000\}$.

Let
\[
 P_4(N)=\frac13+\frac4{135N}-\frac8{2835N^2}
                    -\frac{16}{8505N^3}.
\]
Applying \eqref{eq:poincare} with $R=6$ gives
\begin{equation}
 N^4\bigl(\theta(N)-P_4(N)\bigr)
 =c_4+\frac{c_5}{N}+O(N^{-2}),
 \label{eq:scalederror}
\end{equation}
where
\[
 c_4=\frac{8992}{12629925}\approx0.000711959888914621.
\]
Table~\ref{tab:numerical} illustrates the limit in
\eqref{eq:scalederror}. The numbers are high-precision computations rounded
for display, not interval-certified error bounds; the analytic proof
above supplies the rigorous asymptotic assertion.

\begin{table}[htbp]
\centering
\begin{tabular}{r r r}
\toprule
$N$ & $N^4(\theta(N)-P_4(N))$ & Ratio to $c_4$\\
\midrule
10   & 0.000774546454009 & 1.087907431400\\
20   & 0.000744626622858 & 1.045882829148\\
50   & 0.000725333573049 & 1.018784322464\\
100  & 0.000718695744320 & 1.009461004069\\
200  & 0.000715339856847 & 1.004747413422\\
500  & 0.000713314737917 & 1.001902985018\\
1000 & 0.000712637788046 & 1.000952159162\\
\bottomrule
\end{tabular}
\caption{A numerical illustration of the first omitted coefficient after
four terms. Full working-precision output for five truncation orders is
included in \texttt{artifacts/numerical\_checks.csv}.}
\label{tab:numerical}
\end{table}

The optional numerical run is reproduced with
\begin{lstlisting}[language=bash]
python3 -m pip install -r requirements-optional.txt
python3 code/numerical_check.py
\end{lstlisting}
The README describes the remaining files and the PDF build command.

\section{Conclusion}
\label{sec:conclusion}

Both formulas in A260306 are proved for every $n\geq0$. The analytic part
identifies the asymptotic coefficients uniquely and proves an
$O_R(N^{-R})$ remainder after any fixed number of terms. The algebraic
part is finite: a negative-power identity produces the factorial
$(3n+2)!$ and the denominator $n+i+1$; the exponential generating function
of Stirling numbers gives the double sum; and inclusion--exclusion gives
the triple sum.

The main structural answer is the extraction \eqref{eq:key}, or
equivalently the residue \eqref{eq:residue}. The more elaborate OEIS
expressions are exact finite expansions of that simpler object. The
recurrence, the associated Stirling-number formula, and the computational
checks are consequences and useful cross-checks, not additional
assumptions in the proof.

\clearpage
\appendix
\section{Minimal exact coefficient generator}
\label{app:code}

The following standalone routine computes all coefficients from $c_0$
through $c_K$ using \eqref{eq:hrec} and \eqref{eq:ch}. It uses no symbolic
algebra package and performs no floating-point operations.

\begin{lstlisting}
from fractions import Fraction
from math import factorial


def ramanujan_coefficients(K: int) -> list[Fraction]:
    if not isinstance(K, int) or isinstance(K, bool) or K < 0:
        raise ValueError("K must be a nonnegative integer")
    h = [Fraction(1)]
    for m in range(1, 2 * K + 2):
        convolution = sum(
            (h[m-j] * h[j] / (j+1) for j in range(1, m)),
            Fraction(0),
        )
        h.append(
            (h[m-1] / m - convolution) * Fraction(m+1, m+2)
        )
    return [
        Fraction(int(n == 0)) - 2**n * factorial(n) * h[2*n+1]
        for n in range(K + 1)
    ]


for n, coefficient in enumerate(ramanujan_coefficients(10)):
    print(n, coefficient.numerator, coefficient.denominator)
\end{lstlisting}

The fuller \texttt{code/verify.py} in the archive includes the independent
algorithms and validation checks discussed in
Section~\ref{sec:verification}. The exact checks were executed with
Python 3.13.5; the supplied scripts use syntax supported by Python 3.10
and later.

\begin{thebibliography}{9}

\bibitem{oeis260306}
The OEIS Foundation Inc.,
\emph{A260306: Numerators in Ramanujan's asymptotic expansion of theta(n)},
entry authored by Vladimir Reshetnikov, November 10, 2015, with subsequent
contributions.
\url{https://oeis.org/A260306}.
Text record consulted September 20, 2026:
\url{https://oeis.org/search?q=id:A260306&fmt=text}.

\bibitem{oeis065973}
The OEIS Foundation Inc.,
\emph{A065973: Denominators in an asymptotic expansion of Ramanujan}.
\url{https://oeis.org/A065973}.
Consulted September 20, 2026.

\bibitem{osullivan}
Cormac O'Sullivan,
\emph{Ramanujan's approximation to the exponential function and
generalizations}, arXiv:2205.08504 [math.NT], 2022.
See the introduction and the discussion of associated Stirling numbers
in Section 7.
\url{https://arxiv.org/abs/2205.08504}.

\end{thebibliography}
\end{document}
